% Unidad U015. Acoplamiento de horquillas racionales, refinamiento TPK y
% certificación a profundidad arbitraria.

\label{gfp:b:chap:precision-arbitraria-constantes}

Esta unidad tiene estatuto de certificación posterior. Los prefijos y sus
extensiones compatibles proceden de la relación interna
\(\operatorname{Ext}\), de la supervivencia coinductiva y de las leyes de
transición APP--TRIT--TPK. Las horquillas racionales que siguen no eligen
entre los \(729\) elementos de una fibra: localizan en la carta arquimediana
la extensión ya generada y certifican su precisión a profundidad arbitraria.

\section{Separación tipada de los \(4\) índices de la prolongación nonádica}

Para cada carácter \(\chi\in\{\pi,e,\varphi\}\) intervienen cuatro
parámetros de naturaleza diferente:
\begin{enumerate}
 \item \(n\), orden de la horquilla analítica racional;
 \item \(K\), nivel de refinamiento del Topological Phase Kernel;
 \item \(L_K=6K\), longitud ternaria del prefijo publicado;
 \item \(M\), número de posiciones decimales solicitado.
\end{enumerate}
No se postula \(n=K\), \(n=L_K\), \(K=M\) ni una proporcionalidad fija
entre ellos. El acoplamiento se construye mediante inclusiones de intervalos
decididas por aritmética entera.

Las partes enteras son
\[
 m_\pi=3,
 \qquad
 m_e=2,
 \qquad
 m_\varphi=1,
\]
y los caracteres fraccionarios son
\[
 \rho_\pi=\pi-3,
 \qquad
 \rho_e=e-2,
 \qquad
 \rho_\varphi=\varphi-1.
\]
Las tres familias de horquillas semiabiertas se escriben
\[
 J_{\chi,n}=[\ell_{\chi,n},h_{\chi,n}).
\]
Sus extremos son racionales, contienen a \(\rho_\chi\) y satisfacen
\[
 \ell_{\chi,n}\nearrow\rho_\chi,
 \qquad
 h_{\chi,n}\searrow\rho_\chi,
 \qquad
 h_{\chi,n}-\ell_{\chi,n}\longrightarrow0.
\]
Para \(\pi\), los extremos proceden de las sumas alternadas de la identidad
de Machin demostrada en el \cref{gfp:b:chap:biografia-pi-autonoma}; para \(e\),
de las sumas factoriales y su resto racional del
\cref{gfp:b:chap:biografia-completa-e}; para \(\varphi\), de los cocientes de
Fibonacci y Cassini del \cref{gfp:b:chap:biografia-completa-phi}.

\section{Cilindros ternarios y fibra ambiente inmediata}

Un prefijo ternario de longitud \(L_K=6K\) se representa por un entero
\(N_K\) con
\[
 0\leq N_K<3^{L_K}.
\]
Su cilindro semiabierto es
\begin{equation}
 C(N_K,L_K)=
 \left[
 \frac{N_K}{3^{L_K}},
 \frac{N_K+1}{3^{L_K}}
 \right).
 \label{gfp:b:eq:cilindro-nivel-K}
\end{equation}
La siguiente emisión añade seis trits. La fibra ambiente inmediata contiene
\(3^6=729\) elementos, indexados por \(u\in\{0,\ldots,728\}\):
\begin{equation}
 C_{K,u}=
 \left[
 \frac{729N_K+u}{3^{L_K+6}},
 \frac{729N_K+u+1}{3^{L_K+6}}
 \right).
 \label{gfp:b:eq:729-subcilindros}
\end{equation}
Los \(C_{K,u}\) son disjuntos y
\[
 C(N_K,L_K)=\bigsqcup_{u=0}^{728}C_{K,u}.
\]
El número \(729\) cuenta coordenadas de esta fibra ambiente; no cuenta
semillas, emisiones distintas, órbitas ni estados sobrevivientes.

\subsection{Comparación cruzada de numeradores y denominadores en la horquilla racional}

Sea \(J=[\ell,h)\subset C(N_K,L_K)\), con
\(\ell,h\in\mathbb Q\). Escribimos \(S_K=3^{L_K+6}\) y definimos
\begin{align}
 j_{\min}(J,K)
 &=\max\!\left(729N_K,\left\lfloor S_K\ell\right\rfloor\right),
 \label{gfp:b:eq:jmin-diagonal}\\
 j_{\max}(J,K)
 &=\min\!\left(729N_K+728,\left\lceil S_Kh\right\rceil-1\right).
 \label{gfp:b:eq:jmax-diagonal}
\end{align}

\begin{lemma}[Criterio de localización unitaria]
\label{gfp:b:lem:criterio-localizacion-unitaria}
La horquilla \(J\) está contenida en un único elemento de la fibra
\eqref{gfp:b:eq:729-subcilindros} si y sólo si
\begin{equation}
 \boxed{j_{\min}(J,K)=j_{\max}(J,K).}
 \label{gfp:b:eq:jmin-igual-jmax}
\end{equation}
En ese caso, si \(j\) es el valor común,
\[
 u=j-729N_K,
 \qquad
 N_{K+1}=j.
\]
\end{lemma}

\begin{proof}
Los enteros \(j\) cuyos intervalos elementales de longitud \(S_K^{-1}\)
intersectan \(J\) forman el intervalo entero
\[
 \left[
 \lfloor S_K\ell\rfloor,
 \lceil S_Kh\rceil-1
 \right]\cap
 [729N_K,729N_K+728].
\]
Éste contiene exactamente un elemento si y sólo si sus extremos coinciden.
La igualdad identifica el numerador del subcilindro y, al truncar sus seis
últimos trits, se recupera \(N_K\).
\end{proof}

\section{Certificación diagonal mínima}

La generación no fija de antemano un mismo orden analítico para todos los
niveles. Una vez producida la extensión compatible, el orden de la
horquilla se aumenta sólo hasta certificarla de manera unitaria.

\begin{definition}[Recurrencia diagonal de certificación]
\label{gfp:b:def:recurrencia-diagonal-constantes}
Fijada por \(\operatorname{Ext}\) una familia de cilindros compatibles
\(C(N_K,L_K)\) cuya publicación contiene a \(\rho_\chi\), se elige un orden
inicial \(n_\chi(K_0-1)\geq1\) y definimos recursivamente
\begin{equation}
 n_\chi(K)=
 \min\left\{
 n\geq n_\chi(K-1):
 j_{\min}(J_{\chi,n},K)=j_{\max}(J_{\chi,n},K)
 \right\},
 \label{gfp:b:eq:nchi-K-minimo}
\end{equation}
y verificamos
\begin{equation}
 N_{\chi,K+1}
 =j_{\min}(J_{\chi,n_\chi(K)},K).
 \label{gfp:b:eq:Nchi-K-mas-uno}
\end{equation}
El conjunto
\[
 \Delta_\chi=
 \{(n_\chi(K),K):K\geq K_0\}
\]
es la diagonal analítico--nonádica del carácter \(\chi\).
\end{definition}

\begin{theorem}[Existencia y unicidad de la diagonal de certificación]
\label{gfp:b:thm:existencia-unicidad-diagonal}
Para \(\chi\in\{\pi,e,\varphi\}\), la recurrencia
\eqref{gfp:b:eq:nchi-K-minimo} está definida en todo nivel finito. El índice
\(n_\chi(K)\) y la coordenada \(N_{\chi,K+1}\) son únicos.
\end{theorem}

\begin{proof}
Supongamos generada internamente la extensión cilíndrica de nivel \(K\) que contiene a
\(\rho_\chi\). Como \(\rho_\chi\) es irracional, no pertenece a ninguna
frontera ternaria \(3^{-L}\mathbb Z\). Por ello está en el interior de un
único subcilindro \(C_{K,u_*}\), a distancia positiva de sus dos extremos.
Las horquillas \(J_{\chi,n}\) contienen a \(\rho_\chi\) y su diámetro
tiende a cero; existe, por tanto, un \(n\) para el que la horquilla queda
contenida en \(C_{K,u_*}\). El conjunto del mínimo en
\eqref{gfp:b:eq:nchi-K-minimo} es no vacío y, por el buen orden de
\(\mathbb N\), posee un menor elemento. El
\cref{gfp:b:lem:criterio-localizacion-unitaria} prueba que la horquilla la
localiza de manera unitaria y que el numerador localizado coincide con el ya
generado. La inducción sobre \(K\) completa el argumento.
\end{proof}

\begin{corollary}[Ausencia de parámetros decimales]
La diagonal se calcula con sumas, productos, potencias, divisiones euclídeas,
pisos y techos de enteros. No recibe una cifra decimal de \(\pi\), \(e\) o
\(\varphi\) como dato de certificación.
\end{corollary}

\section{Compatibilidad con la conexión nonádica conjunta}

Sea \(\widehat x_{\chi,K}\) el estado enriquecido del canal \(\chi\) en
el nivel \(K\). Sus coordenadas incluyen el prefijo, el cilindro, la fase,
el cociente, el acarreo, la memoria vertical no acotada, el término terminal
y el registro de procedencia. La fase residual es
\[
 g_0(K)=[K]_9\in\mathbb Z/9\mathbb Z,
\]
mientras la etiqueta pública es
\[
 g(K)=1+((K-1)\bmod9)\in\{1,\ldots,9\}.
\]
La memoria vertical \(q_{\rm mem}(K)\in\mathbb Z_{\geq0}\) no se reduce
módulo \(12\); sólo su proyección dodecafásica
\[
 \beta(K)=q_{\rm mem}(K)\bmod12
\]
lo hace.

Para cada \(K\), sea
\[
 \Gamma_{9,K}:\widehat X_K\longrightarrow\widehat X_{K+9}
\]
la holonomía de una vuelta de fase. Entonces
\[
 g_0(K+9)=g_0(K),
 \qquad
 q_{\rm mem}(K+9)=q_{\rm mem}(K)+1,
 \qquad
 \beta(K+9)=\beta(K)+1\pmod{12}.
\]

\begin{theorem}[Retorno de fase con conservación genealógica]
\label{gfp:b:thm:retorno-fase-conservacion-genealogica}
La igualdad de fase no reinicia el estado:
\[
 g_0(K+9)=g_0(K),
 \qquad
 \widehat x_{\chi,K+9}\ne\widehat x_{\chi,K}.
\]
Las profundidades \(27\), \(54\) y \(108\) son las composiciones tipadas
\begin{align*}
 \Gamma_{27,K}
 &=\Gamma_{9,K+18}\circ\Gamma_{9,K+9}\circ\Gamma_{9,K},\\
 \Gamma_{54,K}
 &=\Gamma_{9,K+45}\circ\cdots\circ\Gamma_{9,K},\\
 \Gamma_{108,K}
 &=\Gamma_{9,K+99}\circ\cdots\circ\Gamma_{9,K}.
\end{align*}
En ningún caso son reinicios del prefijo, del acarreo, del terminal o de la
memoria.
\end{theorem}

\begin{proof}
La primera coordenada se calcula en \(\mathbb Z/9\mathbb Z\), mientras
\(q_{\rm mem}\) pertenece a un monoide no acotado. Cada aplicación de
\(\Gamma_{9,K}\) añade una unidad de memoria, seis bloques por nivel
intermedio y la composición correspondiente de cociclos. Por ello la fase
puede coincidir aunque el estado enriquecido no lo haga. Las fórmulas de
profundidad muestran los índices de cada mapa y evitan componer como si todos
fueran endomorfismos del mismo espacio.
\end{proof}

\subsection{Compresión nonádica, evaluación arquimediana y término terminal del cilindro}

La realización operatorial de una vuelta satisface
\begin{equation}
 U_\Gamma J=JT+\eta,
 \qquad
 J^*\eta=0,
 \qquad
 I-T^*T=\eta^*\eta.
 \label{gfp:b:eq:compresion-expresion-precision}
\end{equation}
La contracción visible estricta se reserva para
\[
 M_9=\frac59V_9;
\]
no se identifica con el operador de compresión \(T\). La telescopía de
profundidad \(L\) conserva el término terminal:
\begin{equation}
 S_0=
 \sum_{r=0}^{L-1}M_9^{*r}D_9^*D_9M_9^r
 +M_9^{*L}S_0M_9^L.
 \label{gfp:b:eq:telescopia-terminal-precision}
\end{equation}
Omitir el último sumando antes de probar su desaparición borraría
información de frontera y convertiría el retorno de fase en un reinicio
falso.

\section{Sistema inverso de historias compatibles}

Sea \(\mathcal H_K^\chi\) el conjunto de historias enriquecidas admisibles
del canal \(\chi\) a profundidad \(K\), y sea
\[
 p_{N,K}:\mathcal H_N^\chi\longrightarrow\mathcal H_K^\chi
 \qquad(N\geq K)
\]
el truncamiento. Definimos
\[
 \operatorname{Surv}_K^{(N)}(\chi)
 =p_{N,K}(\mathcal H_N^\chi)
\]
y
\begin{equation}
 \operatorname{Surv}_K(\chi)
 =\bigcap_{N\geq K}\operatorname{Surv}_K^{(N)}(\chi).
 \label{gfp:b:eq:supervivencia-todas-profundidades}
\end{equation}
El objeto coinductivo es
\begin{equation}
 X_\chi=
 \varprojlim_K
 \bigl(\operatorname{Surv}_K(\chi),p_{K+1,K}\bigr).
 \label{gfp:b:eq:limite-inverso-caracteres}
\end{equation}

\begin{theorem}[Sección compatible bajo prolongabilidad enriquecida]
\label{gfp:b:thm:seccion-compatible-caracteres}
Supóngase además que, sobre cada subcilindro numérico generado, la
relación \(\operatorname{Ext}\) es no vacía, univaluada y compatible con los
truncamientos. Entonces \(\operatorname{Ext}\) y la supervivencia producen una familia
\[
 x_\chi=(\widehat x_{\chi,K})_{K\geq K_0}
\]
que satisface
\[
 p_{K+1,K}(\widehat x_{\chi,K+1})
 =\widehat x_{\chi,K}.
\]
Por tanto, \(x_\chi\in X_\chi\).
\end{theorem}

\begin{proof}
El subcilindro de nivel \(K+1\) tiene numerador
\(N_{\chi,K+1}=729N_{\chi,K}+u\). El truncamiento elimina sus seis últimos
trits y recupera \(N_{\chi,K}\). Las demás coordenadas se obtienen por las
leyes functoriales de acarreo, memoria, terminal y fase. La hipótesis de
univaluación de \(\operatorname{Ext}\) fija el subcilindro y la extensión
enriquecida; la unicidad local del
\cref{gfp:b:lem:criterio-localizacion-unitaria} certifica después su
localización arquimediana, y la compatibilidad de \(\operatorname{Ext}\)
proporciona la igualdad de truncamiento.
\end{proof}

\section{Celdas decimales y precisión prescrita}

Para \(M\geq1\) y \(q\in\{0,\ldots,10^M-1\}\), sea
\[
 D_{M,q}=
 \left[
 \frac{q}{10^M},
 \frac{q+1}{10^M}
 \right).
\]
Una horquilla \(J=[\ell,h)\) está contenida en una única celda decimal si
y sólo si
\begin{equation}
 \boxed{
 \left\lfloor10^M\ell\right\rfloor
 =\left\lceil10^Mh\right\rceil-1=q.}
 \label{gfp:b:eq:criterio-celda-decimal}
\end{equation}

\begin{definition}[Nivel mínimo de certificación decimal]
Para \(\chi\in\{\pi,e,\varphi\}\), definimos \(K_\chi(M)\) como el menor
nivel para el que existen un orden analítico \(n\) y un entero \(q\) tales
que
\begin{enumerate}
 \item \(J_{\chi,n}\) satisface
       \eqref{gfp:b:eq:criterio-celda-decimal};
 \item \(J_{\chi,n}\) queda contenido en el cilindro generado por el
       estado \(\widehat x_{\chi,K}\);
 \item la transición al nivel siguiente satisface
       \eqref{gfp:b:eq:jmin-igual-jmax}.
\end{enumerate}
El par \((K,n)\) se minimiza lexicográficamente.
\end{definition}

\begin{theorem}[Certificación a precisión decimal arbitraria de las salidas coinductivas]
\label{gfp:b:thm:precision-decimal-arbitraria}
Para cada \(\chi\in\{\pi,e,\varphi\}\) y cada \(M\geq1\), existe
\(K_\chi(M)<\infty\). La sección generada determina un único prefijo decimal
de \(M\) posiciones, y la horquilla reconoce y certifica ese prefijo. Si
\(M'<M\), el prefijo de orden \(M\) trunca al de
orden \(M'\), y las celdas correspondientes están encajadas.
\end{theorem}

\begin{proof}
Los tres caracteres son irracionales: \(\varphi\) lo es por su polinomio
cuadrático con discriminante no cuadrado; la irracionalidad de \(e\) y
\(\pi\) se toma de sus teoremas clásicos de caracterización, posteriores a
la construcción finita. Ninguno pertenece, pues, a una frontera decimal o
ternaria racional. Para un \(M\) fijo, la distancia del carácter al conjunto
finito de fronteras decimales relevantes es positiva. Las horquillas
racionales contraen su diámetro hasta quedar en una única celda decimal.
El \cref{gfp:b:thm:existencia-unicidad-diagonal} acopla esa horquilla con un
cilindro TPK y con una extensión única. El mínimo lexicográfico existe por
buen orden. La compatibilidad de la sección prueba el encajamiento al
aumentar \(M\).
\end{proof}

\begin{corollary}[Unicidad real]
Las celdas encajadas tienen diámetros \(10^{-M}\to0\). Su intersección es
un singleton, identificado por las caracterizaciones precedentes con
\(\pi\), \(e\) o \(\varphi\), respectivamente.
\end{corollary}

\section{Equivalencia de \(2\) realizaciones finitas de la ley uniforme}

Los testigos de \(1,000\) y \(10,000\) posiciones no definen los
caracteres ni limitan la prolongación. Son ejecuciones de la misma recurrencia
uniforme con dos profundidades distintas.

\begin{table}[htbp]
\centering
\caption{Dimensiones exactas de las ejecuciones certificadas, por canal.}
\label{gfp:b:tab:dimensiones-ejecuciones-1000-10000}
\begin{tabular}{lrr}
\toprule
Objeto contado & \(1,000\) posiciones & \(10,000\) posiciones\\
\midrule
Particiones exhaustivas de la fibra ambiente & 396 & 3501\\
Trits transportados & 2376 & 21006\\
Pares de igual fase \((K,K+9)\) & 387 & 3492\\
Vueltas nonádicas completas & 43 & 388\\
Elementos examinados por partición & 729 & 729\\
Extensiones compatibles conservadas por partición & 1 & 1\\
\bottomrule
\end{tabular}
\end{table}

Para \(10,000\) posiciones,
\[
 B=3501=389\cdot9,
 \qquad
 6B=21006,
 \qquad
 3^{6B}>10^{10000}.
\]
Los primeros \(396\) bloques de la ejecución larga coinciden con la
ejecución de \(1,000\) posiciones. En cada uno de los \(3492\) pares de
igual fase se preserva la coordenada de fase y se prolongan el prefijo, la
profundidad, el cilindro y la memoria. El certificado registra también las
repeticiones eventuales de proyecciones locales sin confundirlas con retorno
del estado completo.

\section{Separación entre generación y verificación}

\begin{definition}[Generador estructural]
El generador recibe el catálogo finito, \(L_0,L_1\), el estado inicial, las
leyes de conexión y la relación interna \(\operatorname{Ext}\). En cada nivel:
\begin{enumerate}
 \item enumera los \(729\) elementos de la fibra ambiente;
 \item aplica \(\operatorname{Ext}\) y conserva la única extensión compatible;
 \item actualiza fase, cociente, acarreo, memoria y terminal;
 \item escribe el bloque ternario, el recibo de selección y el registro genealógico.
\end{enumerate}
\end{definition}

\begin{definition}[Verificador posterior]
El verificador recibe las salidas cerradas del generador y las tres familias
de horquillas racionales; calcula \(j_{\min}\) y \(j_{\max}\), certifica su
igualdad con la extensión ya producida, reconstruye las
particiones, comprueba los cuadrados de truncamiento, los pares
\((K,K+9)\), las sumas de partición
\[
 \operatorname{killed}_{\mathrm{izq}}+1+\operatorname{killed}_{\mathrm{der}}=729,
\]
las identidades de contenido y, al final, la coincidencia con expansiones de referencia.
\end{definition}

\begin{theorem}[Aislamiento de los datos objetivo]
\label{gfp:b:thm:aislamiento-datos-objetivo-precision}
Las expansiones decimales de referencia y los datos metrológicos no
intervienen en la selección de una extensión compatible. La comparación se
ejecuta después de escribir los tres prefijos.
\end{theorem}

\begin{proof}
La entrada matemática de cada decisión generativa consiste en el catálogo
finito, \(L_0,L_1\), el estado heredado y la relación interna
\(\operatorname{Ext}\). Los extremos racionales, potencias de \(3\), pisos y
techos pertenecen al verificador posterior. El generador termina antes de que
éste abra las horquillas o los controles decimales. La dirección del flujo de
datos es
\[
 \text{estructura}\longrightarrow
 \text{prefijos y ledgers}\longrightarrow
 \text{verificación posterior};
\]
no existe una flecha en sentido contrario.
\end{proof}

\section{Realización finita del certificado de precisión arbitraria}

Una realización material del procedimiento produce, para cada uno de los tres caracteres, un prefijo de \(10^4\) posiciones. Este corte finito es un control posterior de la construcción y no interviene en la elección de cilindros, acarreos ni extensiones compatibles. La proposición demostrada conserva alcance arbitrario en \(M\); la precisión \(10^4\) certifica una ejecución concreta, pero no limita el generador coinductivo.

\section{Obstrucciones diferenciadas de la generación y de su certificación a precisión arbitraria}

\begin{remark}[Certificado y control de validez: Refutación de la generación a precisión arbitraria]
Los fallos de órbita, extensión, truncamiento o retorno refutan la generación;
los fallos de horquilla, orden o control decimal refutan el certificado
analítico correspondiente; y una discrepancia entre ambos refuta la
identificación conjunta de la salida publicada. Con esta separación de tipos,
la afirmación completa queda refutada si:
\begin{enumerate}
 \item se identifica cualquiera de los índices \(n,K,L_K,M\) sin un
       teorema de enlace;
 \item una horquilla deja de contener su carácter o no contrae su diámetro;
 \item una etapa satisface \(j_{\min}\ne j_{\max}\) después del orden
       declarado como mínimo;
 \item una extensión no trunca al cilindro del nivel anterior;
 \item el retorno de fase reinicia prefijo, acarreo, terminal o memoria;
 \item se omite el término terminal en
       \eqref{gfp:b:eq:telescopia-terminal-precision} sin demostrar su anulación;
 \item los primeros \(396\) bloques de la ejecución larga difieren de la
       ejecución corta;
 \item un control decimal o metrológico es leído antes de cerrar la salida;
 \item una cifra certificada difiere de la salida material del generador;
 \item la prueba depende de que la precisión solicitada sea \(1,000\) o
       \(10,000\), en vez de un \(M\) arbitrario.
\end{enumerate}
\end{remark}
